\section{Introduction}
\label{sec:introduction}

Methods built on random convolutional kernels combine accuracy and speed better than most
other approaches to time series classification
\citep{dempster2020rocket,dempster2021minirocket,tan2022multirocket,middlehurst2024bakeoff}.
Hydra \citep{dempster2023hydra} organises such kernels into groups that compete at every
time step and counts how often each kernel wins, which links convolutional methods with
dictionary methods that count patterns. Together with MultiRocket features, Hydra is among
the most accurate methods on the UCR archive \citep{middlehurst2024bakeoff}.

The competition has a cost. To find the winner of a group of $k$ kernels of length $\ell$,
Hydra evaluates every kernel: $k\ell$ multiplications per group and time step. Work on edge
devices makes each kernel cheaper. NanoHydra \citep{cioflan2025nanohydra} uses binary
kernels and integer arithmetic on a microcontroller, but still evaluates every kernel of a
group, and in its own evaluation it is 1.5 points less accurate than Hydra on the UCR
archive (83.81\% against 85.31\%), although its numbers are the best of several runs per
dataset. This paper asks a different question: can the winner of a group be found without
evaluating its kernels one by one, without losing accuracy?

We answer it with groups that have a structure. Each group of \hh{} is the vertex set of a
Hanner polytope, a convex body built from segments by products and $\ell_1$ sums
\citep{hanner1956intersections,hansen1981structure}. In the canonical coordinates, every
vertex is a ternary vector in $\{-1,0,1\}^\ell$, and the polytope's support function follows
the tree of its construction. The winning vertex of a group with up to $2^\ell$ vertices
then costs $\ell$ absolute values and $\ell-1$ additions or comparisons, and no
multiplications. The winners are sign patterns on subsets of the taps, the codes of
hyperplane and cross-polytope hashing \citep{charikar2002simhash,andoni2015crosspolytope}.
Hanner polytopes also have an extremal property: they minimise the volume product among
unconditional bodies, and an unrefereed 2026 manuscript claims that they are the equality
cases of the symmetric Mahler conjecture in general \citep{openai2026mahler}. The method does
not depend on this property (Section~\ref{sec:background}).

A cheaper group is only useful if it does not cost accuracy, and a fast new kernel family
invites the question of which of its properties matter. We therefore tested a
non-inferiority hypothesis that we fixed in advance, with two controls matched to the
Hanner groups in size and sparsity: random ternary groups, and random ternary groups that
are centrally symmetric like every Hanner polytope.

Our contributions are:
\begin{itemize}
\item the Hanner kernel family, with a proof that the winner of a group follows its tree in
  $O(\ell)$ operations, including an exact rule for ties (Section~\ref{sec:method});
\item two theory-matched controls that separate the tree structure from ternary weights
  and from central symmetry (Section~\ref{sec:controls});
\item a pre-registered evaluation on 138 UCR and UEA datasets with 30 resamples each, a
  non-inferiority decision rule and a controlled single-core timing protocol
  (Section~\ref{sec:protocol});
\item a NumPy reference implementation, Numba kernels that reproduce it bit for bit,
  aeon-compatible estimators, and the full per-case results (Section~\ref{sec:reproducibility}).
\end{itemize}

The pre-registered hypothesis is accepted (Section~\ref{sec:results}). \hh{} is practically
equivalent to Hydra in accuracy, $-0.06$ points with a posterior probability of 0.988 of
being not worse, and to the matched controls. Its single-core transform is $11.2\times$
faster than Hydra's, and it needs $23.8\times$ less memory. The result is one of
equivalence, not improvement, and it has limits: on the UEA datasets alone, \emph{not
worse} is not shown, the gap to Hydra on balanced accuracy and macro-F1 is inconclusive,
and MultiRocket and HIVE-COTE~2.0 remain more accurate. The tree structure contributes
speed, not accuracy. Section~\ref{sec:background} reviews Hydra and Hanner polytopes,
Section~\ref{sec:related} discusses related work, and Section~\ref{sec:discussion} the
limitations.
